\section{Conclusions}
\label{iv:conclusiones}
\begingroup
\setlength{\parskip}{0pt}
\setlength{\abovedisplayskip}{5pt}
\setlength{\belowdisplayskip}{5pt}

La structure discrète du continu développée dans cet article
soutient une réalisation exceptionnelle qui atteint \(V^\natural\) et
une réalisation opératorielle qui conserve la dualité compacte.
Le lien est construit : l’état APP--TRIT--TPK transmet
l’orientation, les feuillets, les résidus, les quotients, le bord et la mémoire à ses
coordonnées numériques, au registre d’incidence et aux représentations
de phase et d’action. La contribution centrale consiste à relier ces
réalisations au moyen de codes, de recollements, de produits de vertex,
de projecteurs et d’opérateurs d’entrelacement à domaines déterminés.

La construction conjointe du continu fournit l’antécédent de
ces applications. Les mêmes prolongements produisent la mesure des
histoires, l’identité de Gram par feuillets avec son résidu terminal,
le défaut d’annulation, l’incidence rectangulaire et le transport
naturel du complexe de mémoire. L’assemblage conserve toutes les
continuations et satisfait \(D^2=0\) sur ses générateurs ; sa droite
déterminante conserve les degrés et les bases. Les poids compatibles,
l’isométrie du transport, le résidu de survie et les conditions
d’une lecture déterminante scalaire maintiennent la fonction concrète que
les preuves leur attribuent. Ces opérations corrélées constituent
la structure commune antérieure aux représentations numériques,
exceptionnelles et dimensionnelles.

% BEGIN R27_FRAMING_CONCLUSION_ES
L’articulation arithmétique et d’incidence est précisée par des résultats
à domaines explicites. Les lois du défaut distinguent le raffinement
du module, la répétition périodique et l’augmentation de dimension. L’opérateur
agrégé conserve un sous-réseau entier d’indice deux. Dans son bloc
quadratique \(3H\), l’unité normalisée \(H\) a pour norme un ; la
récurrence de Pell détermine \(H^n\) et le bloc d’origine évolue
comme \(3^nH^n\).
La bijection entre paires de trits et chiffres nonadiques conserve la retenue,
tandis que la projection du cube ambiant a neuf éléments par fibre.
L’alphabet de \(42\) blocs, les raccordements régionaux et les corrélations
de registres ordonnés rendent visible l’information antérieure aux
représentations linéaires.

% BEGIN R28_FRAMING_CONCLUSION_ES
Les cinq clôtures régionales atteignent le minimum huit dans le problème fini de couverture des quatre arêtes prescrites. Le graphe conserve les étiquettes des émissions et le quotient conserve leurs incidences entre classes de phase. L'inversion orientée rend impair le lecteur \(\nu_{120}\) et maintient pair \(\nu_{270}\) ; ces identités explicitent la réponse de la mémoire au changement d'orientation.
% END R28_FRAMING_CONCLUSION_ES

% BEGIN S27_FRAMING
La comparaison régionale détermine trois configurations polarisées :
les étoiles de \(B_K\), \(B_\varphi\) et \(B_{\rm APP}\) ont,
respectivement, les signatures \((3,1,0)\), \((1,0,3)\) et \((1,1,2)\).
La couronne \(\{i:K_i\geq729\}=B_K=H_e\cap P_\pi\) relie
le registre numérique à un quadruplet d’incidence, et la cochaîne
signée fixe \(H^-_\alpha=B_K\sqcup\{5,7\}\).
Le recensement complet situe la signature \((3,1,0)\) dans quinze quadruplets.
La rigidification correspond au repère ordonné
\((H^-_\alpha,H_e,H_\varphi,H_{\rm APP})\) : son stabilisateur
dans \(G_W\) est l’identité. Le résultat fixe la liberté de
permutation des douze positions ; les valeurs de \(K\) conservent
leur détermination par le registre entier et ses opérations de
reconstruction.

Le relèvement entier de la signature régionale ajoute un résultat
complémentaire. Les identités \(p+e+f=q+3c\) et
\(p+e-f=a+3h\) conservent les quotients qui accompagnent les résidus,
et le cocycle de retenue rend leur composition compatible avec
l’associativité de la somme. Dans le tableau complet à \(27\) colonnes,
les triplets \((0,0,0)\) et \((1,2,0)\) ont \(q=a=0\), tandis que
leurs paires \((c,h)\) sont \((0,0)\) et \((1,1)\).
Chacune des réductions \(\bar c,\bar h\), ajoutée séparément
à l’évaluation de \(1,p,e,f\), augmente le rang de quatre à cinq.
La lecture complète \((q,a,c,h)\) conserve à son tour des fibres :
\((1,2,0)\) et \((2,1,0)\) produisent la même donnée
\((0,0,1,1)\). Les identités de transport et les bornes du cylindre
précisent donc l’information supplémentaire de la mémoire entière
et la multiplicité des antécédents qui demeure dans chaque lecture.
Cette articulation conserve conjointement le contenu numérique et
l’incidence de l’histoire régionale.

% END S27_FRAMING


L’identité \(C^2=5I_6\) provient de la corrélation quadratique évaluée
sur les cinq positions finies. La saturation et la neutralité duale laissent
une paire de bords dont le représentant dépend de l’orientation
déclarée. Les contrôles de clôture distinguent cette sélection des
variations de bord et de position du registre. 

Le résultat exceptionnel s’obtient en suivant l’incidence de \(K\)
jusqu’à l’algèbre de vertex. Le registre dodécaphasique est construit avant
sa lecture rationnelle et la sélection des supports. Sa configuration
signée intervient dans le drapeau de Witt ; le code ternaire détermine
le recollement de \(A_2^{12}\), et le drapeau avec la métrique, l’origine et la
chambre déclarées détermine le voisin de Leech. Les extensions
d’ordres deux et trois de son algèbre réticulaire réalisent \(V^\natural\)
par application des théorèmes de construction correspondants, dont
les hypothèses sont vérifiées dans le corridor réticulaire. Les isomorphismes
choisis transportent le vide, le vecteur conforme, la graduation, le produit et les
traces ; ils n’identifient pas un automorphisme d’ordre deux à un autre d’ordre
trois. Le Monstre agit dans \(V^\natural\), et Moonshine détermine ses
traces graduées. La lecture rationnelle du registre et les chiffres des
constantes conservent d’autres domaines : leur connexion à l’algèbre est la
chaîne d’objets et d’applications construite, non une action sur leurs chiffres.

Dans la réalisation
exceptionnelle, l’identité rationnelle démontre
\(J-(t_3+12)\in3^9q\mathbb Z[[q]]\) ; son premier coefficient détermine
l’optimalité de l’exposant uniforme. La preuve formelle conserve sa
portée à tous les degrés et les vérifications tronquées conservent
leur fonction de comparaison finie.
% END R27_FRAMING_CONCLUSION_ES

Le même \(K\) détermine une polarisation géométrique et une référence
de récupération. Dans la représentation réelle fixée, le vecteur non nul
\(u_K\) détermine une droite dans \(V_{11}\) ; sa norme quadratique
exacte normalise le terme de rang un du projecteur orthogonal de rang
dix, dont la loi de covariance est spécifiée.
L’orbite dodécaphasique complète est inversible et permet d’identifier
des opérateurs à partir de leurs réponses avec des bornes exactes de stabilité.
Son codage rationnel conserve les douze blocs avec le système de coordonnées fixé ;
le bilan bilatéral récupère l’état à partir de deux canaux définis par des isométries,
et son extension unitaire récupère conjointement l’état et la mémoire
entrante. L’invariance d’autres observables conserve la condition
de commutation prouvée. La fonction arithmético-géométrique de \(K\)
comprend donc l’incidence, la direction et l’identification, en plus
de son évaluation scalaire.

Le prolongement unifère de la représentation de Weyl établit une
réalisation continue de la dimension traciale. Les opérateurs de phase
et de décalage engendrent chaque algèbre matricielle finie, et les
inclusions de fibre complète conservent l’unité et la trace normalisée.
La clôture uniformément hyperfinie nonadique a une représentation
traciale dont la clôture faible est un facteur hyperfini de type \(II_1\) ;
en elle, les rangs normalisés atteignent toutes les dimensions traciales
de l’intervalle unité. La continuation marquée utilise une autre inclusion
et produit une autre clôture. Le résultat identifie l’opération de
raffinement responsable de la conservation de la normalisation globale.

\Needspace{14\baselineskip}
L’action et la paire angulaire déterminent la forme compacte, au lieu de
recevoir un rayon prescrit. Le renouvellement régional normalisé et les
sections positives d’action fixent l’échelle commune de l’ellipse ;
ses coordonnées angulaires fixent le quotient des demi-axes. L’unique rayon
positif satisfait
\[
 R_\star^4=\frac{A-C^*}{A+C^*}
 =\frac{499\alpha-\eta_{\mathrm{ret}}}{501\alpha+\eta_{\mathrm{ret}}}.
\]
L’appartenance au système de coordonnées elliptique et le choix de la racine positive
sont démontrés. Le produit des demi-axes conserve l’action ;
leur échange transforme le rayon en \(R_\star^{-1}\).
Les systèmes de coordonnées angulaires, l’orientation des axes et l’échelle d’action
maintiennent ainsi leurs fonctions respectives dans la forme quadratique.

La dualité T établie est une équivalence scalaire et opératorielle.
Une même involution conserve la forme compacte et conjugue unitairement
les Hamiltoniens autoadjoints de rayons réciproques, y compris leurs
domaines et ceux de leurs formes quadratiques. La covariance de la section
résiduelle exige de transporter \(L_9\) ; sur le quotient correspondant,
l’énergie minimale est bien définie et conserve la dualité.
La normalisation de corde transporte ces identités aux rayons
dimensionnels sans modifier le rayon de forme engendré. La collection
elliptique contient les deux fibres réciproques et conserve la restriction
des domaines et de l’isomorphisme des écrans. La semi-conjugaison
modulaire sur l’axe imaginaire relie leurs paramètres et leurs fonctions,
en conservant le domaine propre des opérateurs de chaque réalisation.

Les écrans transforment la hiérarchie dimensionnelle en applications
spécifiées. Les réductions \(12\to11\to10\) et \(26\to24\)
proviennent, respectivement, de l’élimination du mode uniforme et de la
direction \(u_K\), et de la restriction orthogonale suivie du quotient
isotrope du réseau lorentzien. Le morphisme conjoint identifie
les présentations quotientées et entrelace le secteur compact.
Le graphe conserve la coordonnée dans \(V_{11}\) avec sa projection ;
chaque projection isolée maintient son noyau. Le poinçonnage ternaire
ajoute l’écran combinatoire de longueur onze avec distance cinq
et partition parfaite. Les rangs sont déterminés par ces
opérations, non par une identification entre espaces de même
dimension.

Le prolongement vers la théorie M a un support algébrique
déterminé et une condition dynamique précise. Les changements de direction
et de type arithmétique des trajectoires, composés avec leur donnée
précédente de bord, forment un cocycle exactement additif.
La surface d’univers, l’action de Polyakov et le codomaine des champs
à onze dimensions reçoivent ensuite cette structure au moyen des applications
de réalisation. Le théorème de transport conserve \(Q^2=H\) dans
l’image de l’isométrie sous les entrelacements et domaines déclarés.
L’identification physique complète inclut les champs, les tensions
et la normalisation d’action, le secteur des oscillateurs et ses amplitudes,
le contrôle des anomalies et la limite gravitationnelle de basse énergie de
chaque compactification
(section~\ref{iv:condiciones-realizacion-fisica}).
Ces exigences fixent la dynamique dans laquelle les
structures démontrées sont transportées.

Une chaîne commune est établie, qui relie la construction discrète
du continu à l’algèbre Moonshine, à la polarisation dodécaphasique,
à la réalisation traciale, à la dualité compacte et à la composition d’action.
Ses applications déterminent quelle information demeure, quel quotient est pris
et quelles équations sont conservées. C’est le contenu mathématique que
reçoivent les réalisations dimensionnelles et sur lequel se poursuivent les
développements de statistique, de rayonnement et de structure des espèces.

\par\endgroup
